% language=us

\environment context-2022-style

\startdocument
  [title={The LMTX distribution},
   banner={keep it simple},
   location={context\enspace {\bf 2022}\enspace meeting}]

\starttitle[title=The challenge]

\startitemize
    \startitem
        How can we provide users with a long term stable \LUAMETATEX\ engine that
        doesn't get messed up by uncontrolled patches?
    \stopitem
    \startitem
        How can we ensure a stable current \CONTEXT\ distribution that works with the
        engine as it is when that current was published?
    \stopitem
    \startitem
        How can users get independence of changes in hardware and the operating
        system out of their control?
    \stopitem
    \startitem
        How can we make sure that the right additional resources that we expect to
        be present (with a given version) are kept with the engine and distribution?
    \stopitem
    \startitem
        How can we make installing and running the lot mainly a matter of running the
        single independent binary in a well defined \TEX\ tree?
    \stopitem
    \startitem
        How can we stay stable and stay free from these ever changing fashions,
        pseudo standards, software politics, etc.?
    \stopitem
    \startitem
        In other words how can we best distribute \CONTEXT\ \LMTX\ best?
    \stopitem
\stopitemize

\stoptitle

\starttitle[title=Running]

\startitemize
    \startitem
        We only need one binary: \type {luametatex}.
    \stopitem
    \startitem
        We run \CONTEXT\ with \type {context} which is a link to the binary.
    \stopitem
    \startitem
        We run other scripts with \type {mtxrun} which is a link to the binary.
    \stopitem
    \startitem
        These runners have \type {context.lua} and \type {mtxrun.lua} (with
        suffix) in the same (binary) path.
    \stopitem
    \startitem
        So three files and two links is all what is needed.
    \stopitem
    \startitem
        We assume a standardized setup. One can divert but we're not going to
        waste time on alternatives.
    \stopitem
    \startitem
        Here the resources are all located relative to the binary. This is a
        proven setup (minimals).
    \stopitem
\stopitemize

\stoptitle

\starttitle[title=Directory structure]

This is how one can install it:

\starttabulate[|l|Tl|]
    \NC \notabene {always} \NC /data/context/tex/texmf/ \NC \NR
    \NC \notabene {always} \NC /data/context/tex/texmf-<platform>/bin/ \NC \NR
    \NC \notabene {always} \NC /data/context/tex/texmf-context/ \NC \NR
    \TB
    \NC optional           \NC /data/context/tex/texmf-local/ \NC \NR
    \NC optional           \NC /data/context/tex/texmf-fonts/ \NC \NR
    \NC optional           \NC /data/context/tex/texmf-project/ \NC \NR
\stoptabulate

After running:

\starttyping
mtxrun --generate
mtxrun --script font --reload
\stoptyping

This should work:

\starttyping
context [--autopdf] --global welcome-to-context
\stoptyping

It picks up a file from the tree and runs it on the local path.

\stoptitle


\starttitle[title=A solution]

\startitemize
    \startitem
        Already years ago (and also recently when we met) Mojca and I discussed
        this but we never came to implement it. Of course we blame COVID.
    \stopitem
    \startitem
        It boils down to having several \notabene {repositories} from which we install:
        \startitem
            The \notabene {\CONTEXT\ distribution}, including the source of
            \LUAMETATEX. Basically this is the already present mirror.
        \stopitem
        \startitem
            A repository that has (mostly) \notabene {reliable fonts} that we
            know work as we expect and don't get replaced without them being
            tested.
        \stopitem
        \startitem
            A repository with stable and quality assured \notabene {reliable
            third party modules}.
        \stopitem
        \startitem
            A repository with \notabene {binaries} that match the latest version.
            Users can always generate their own if needed.
        \stopitem
        \startitem
            \notabene {Snapshots} of working combination of these that can get
            installed manually or via an installer.
        \stopitem
    \stopitem
    \startitem
        Basically we're talking of what is (and has always) been the minimals
        but maintained by some more people. New is that we provide the means
        to compile the engine, which is the only program needed.
    \stopitem
    \startitem
        We stay far from dependencies so we have an independent source tree.
        There is the possibility to load libraries, but we provide only very few
        small optional interfaces. Additional ones can be made outside the source
        tree but will not be distributed.
    \stopitem
    \startitem
        For the engine as well as the optional libraries we can set up the
        compile farm but for the later we will not enter a frequent update
        mode: we need little.
    \stopitem
\stopitemize

\stoptitle

\starttitle[title=Compiling]

\startitemize
    \startitem
        Are we ready to distribute the source? Last year we didn't release
        because one important condition was not met. But now we can decide to
        release version 1.00!
    \stopitem
    \startitem
        Although there are plenty of places where I want to go a bit further,
        most is done.
    \stopitem
    \startitem
        Can we expect users to compile binaries? The answer for a complete
        \TEXLIVE\ setup is \quote {no} but for what we have here is can work
        well.
    \stopitem
    \startitem
        In due time we can have recipes on the garden but it should work more
        or less out of the box: \LINUX\ and \MSWINDOWS\ demo.
    \stopitem
    \startitem
        Michal Vlasák made a prototype of setup for compiling some libraries.
    \stopitem
    \startitem
        We need to get \LMTX\ into the next \TEXLIVE. Because we don't rely on
        the build infrastructure Mojca probably has to provide the binaries from
        the farm.
    \stopitem
\stopitemize

\stoptitle

\stoptext
